\chapter{Jeux de teneur de marché}\label{ch:m2:market-making-games}

Un recruteur désigne la fenêtre~: «~faites-moi un \omterm{def:m2:market-making-games:mmam}{make-me-a-market} sur le nombre de fenêtres de cet
immeuble~». On cote 200 à 400. Elle dit~: «~j'achète~». On vient de lui vendre ses fenêtres à 400,
et l'on a appris quelque chose~: elle croit qu'il y en a plus de 400, ou du moins elle est
prête à le parier. La prochaine cotation doit être plus haute. Les sociétés de trading forment et évaluent
des candidats à l'aide de jeux de ce type, parce que le problème de tout teneur de marché s'y trouve~: estimer une
valeur que l'on ne voit pas, coter autour d'elle assez large pour survivre à ceux qui en savent davantage,
assez serré pour traiter avec ceux qui n'en savent pas plus, et apprendre de chaque transaction. Ce
chapitre joue ce jeu au sérieux~: comment faire un \omterm{def:m2:market-making-games:mmam}{make-me-a-market}, comment mettre à jour
son estimation à partir des transactions, comment la sélection adverse fixe la largeur, et ce que disent les sociétés de trading des jeux qu'elles utilisent.

\section{Faire un make-me-a-market}

\begin{definition}[Make-me-a-market, largeur]\label{def:m2:market-making-games:mmam}
Le \emph{make-me-a-market}\index{make-me-a-market} est une \omterm{def:m2:how-bonds-trade:rfq}{demande de cotation}, sur une grandeur incertaine,
d'un cours acheteur auquel on achètera et d'un cours vendeur auquel on vendra, pour une taille donnée,
avant de savoir dans quel sens la contrepartie va traiter. La \emph{largeur}\index{largeur}
du marché est le cours vendeur moins le cours acheteur.
\end{definition}

Le vocabulaire est celui de la corbeille, et le guide public de Jane Street sur ses
entretiens de trading en donne le détail~: «~I'm 2 bid for 10~» propose d'acheter 10 à 2~; «~I have 10 at 4~»
propose de vendre~; «~2 at 4, 10 up~» fait les deux~; l'autre partie dit «~sold~» pour frapper le cours acheteur ou
«~take 'em~» pour lever le cours vendeur, et l'ordre reste en vigueur jusqu'à ce qu'il soit exécuté ou que son auteur dise
«~I'm out~». Le guide conseille, pour un contrat versant le lancer d'un dé, valant 3.5 en
moyenne, consiste à acheter en dessous et à vendre au-dessus de l'espérance, à peser la pire perte
face à son capital, et à arbitrer la probabilité de traiter contre le profit par
transaction~; il suggère «~3 at 4, 10 up~».

Un marché a un centre, c'est-à-dire l'estimation de la valeur par le teneur de marché, et une largeur. Le centre est un
problème de statistique~; la largeur est un problème économique. Coter trop serré revient à offrir à quiconque connaît mieux la valeur
une occasion de traiter contre soi quand on se trompe~; coter trop large, et plus personne ne
traite. Le jeu de cartes de ce chapitre rend les deux points précis.

\begin{definition}[Inventaire]\label{def:m2:market-making-games:inv}
L'\emph{inventaire}\index{inventaire} d'un teneur de marché est la position nette qu'il détient après les
transactions réalisées contre ses cotations, longue après un achat net et courte après une vente nette.
\end{definition}

\section{Mise à jour à partir des transactions}

\begin{definition}[Mise à jour bayésienne]\label{def:m2:market-making-games:bayes}
Une \emph{mise à jour bayésienne}\index{mise à jour bayésienne} révise une loi de probabilité sur une
valeur inconnue après une observation, en proportion de la vraisemblance de cette observation
sous chaque valeur~: a posteriori $\propto$ a priori $\times$ vraisemblance.
\end{definition}

Le contrat du jeu de ce chapitre verse la somme de cinq cartes tirées face cachée dans un
jeu de 52 cartes, les as comptant 1 et les rois 13. Avant toute transaction, la somme vaut 35 en moyenne, avec un
écart-type de 8.03, et elle se situe entre 6 et 64
(la \cref{fig:m2:market-making-games:dist}). À chaque tour arrive un trader~: avec une probabilité
de 0.2, un trader informé, qui connaît la somme et achète si elle est au-dessus du cours vendeur,
vend si elle est sous le cours acheteur, et passe sinon~; sinon arrive un trader non informé, qui traite
dans un sens au hasard, avec une probabilité qui décroît avec la largeur, $e^{-w/8}$, et passe sinon.

\begin{omfigure}
\begin{tikzpicture}
\begin{axis}[width=11cm,height=4.6cm,ybar,bar width=2.2pt,
  xlabel={somme de cinq cartes},ylabel={probabilité},
  tick label style={font=\small},label style={font=\small},
  xmin=4,xmax=66,ymin=0,ymax=0.055,grid=major,grid style={gray!25},scaled y ticks=false,
  yticklabel style={/pgf/number format/fixed,/pgf/number format/precision=2},
  table/col sep=comma]
\addplot[fill=omDef!60,draw=omDef] table[x=sum,y=p]{figdata/markets-2/30-market-making-games/distribution.csv};
\addplot[omThm,very thick,dashed,sharp plot] coordinates {(35,0) (35,0.055)};
\end{axis}
\end{tikzpicture}

\omcaption{Loi de la somme de cinq cartes tirées d'un jeu de 52 cartes, les as comptant 1
et les rois 13~: moyenne 35 (en tirets), écart-type 8.03, de 6 à 64. Valeurs exactes, par
dénombrement. Données~: le tutoriel du chapitre.}\label{fig:m2:market-making-games:dist}
\end{omfigure}

\begin{example}[La première transaction]\label{ex:m2:market-making-games:first}
Le teneur de marché cote 30 à 40, soit une largeur de 10. Un trader non informé traite avec une probabilité
$e^{-10/8} = 0.29$. Si quelqu'un achète à 40, il s'agissait d'un trader informé avec une probabilité de 0.31~;
l'espérance de la somme pour le teneur de marché passe de 35 à 38.13, et sa cotation suivante est
centrée là. Après une vente à 30, elle tombe à 31.87~; après une passe, elle reste à 35, car une
passe est tout aussi probable depuis une somme haute que depuis une somme basse.
\end{example}

Au fil d'une partie, le centre du teneur de marché se rapproche de la vraie somme à mesure que les transactions la révèlent
(la \cref{fig:m2:market-making-games:game})~; la rapidité dépend de la part de traders informés
et de la largeur, car un marché large filtre les transactions informées qui portent
l'information avec les transactions non informées qui paient pour la largeur.

\begin{omfigure}
\begin{tikzpicture}
\begin{axis}[width=11cm,height=4.6cm,
  xlabel={tour},ylabel={milieu du teneur de marché},
  tick label style={font=\small},label style={font=\small},
  xmin=0,xmax=20,ymin=28,ymax=42,grid=major,grid style={gray!25},
  legend columns=2,legend style={font=\footnotesize,at={(0.5,-0.3)},anchor=north,draw=none,fill=none,/tikz/every even column/.append style={column sep=8pt}},
  table/col sep=comma]
\addplot[omDef,very thick,const plot,mark=*,mark size=1.2pt] table[x=round,y=mid]{figdata/markets-2/30-market-making-games/game.csv};
\addplot[omThm,thick,dashed] table[x=round,y=value]{figdata/markets-2/30-market-making-games/game.csv};
\legend{espérance du teneur de marché,somme réelle (31)}
\end{axis}
\end{tikzpicture}

\omcaption{Une partie à une largeur de 11.5~: l'espérance de la somme pour le teneur de marché avant chaque
tour, mise à jour selon la règle de Bayes après chaque transaction et chaque passe, face à la somme réelle, 31. Les
achats précoces font monter l'estimation~; les ventes qui suivent la font redescendre vers la vérité. Le
teneur de marché termine la partie avec un profit de 18.6. Données~: le tutoriel du chapitre,
avec graine fixée.}\label{fig:m2:market-making-games:game}
\end{omfigure}

\section{La sélection adverse autour d'une table}

La définition de la sélection adverse donnée par le guide correspond à celle du livre~1,
chapitre~1~: les transactions auxquelles on accède sont moins avantageuses qu'il n'y paraît, parce que
quelqu'un d'autre sélectionne la transaction contre soi. Sa forme pratique est une question à se poser avant
chaque cotation~: sachant que cet ordre sera exécuté, que vaut réellement l'actif, et voudrait-on encore
réaliser la transaction~? Glosten et Milgrom en ont fait un modèle en 1985~: un spécialiste confronté à des
intervenants informés et à des intervenants non informés fixe son cours acheteur et son cours vendeur à
l'espérance conditionnelle à une vente et à un achat, et la présence de négociateurs mieux informés produit une
fourchette positive, même lorsque le spécialiste est neutre au risque et ne réalise aucun profit espéré.

\begin{definition}[Malédiction du vainqueur]\label{def:m2:market-making-games:curse}
La \emph{malédiction du vainqueur}\index{malédiction du vainqueur} est la tendance du vainqueur d'une enchère
entre enchérisseurs aux estimations bruitées d'une valeur commune à l'avoir surestimée~: remporter
l'enchère est en soi la preuve que son estimation était trop élevée.
\end{definition}

Un teneur de marché est toujours le vainqueur d'une telle enchère~: sa cotation est celle sur laquelle
quelqu'un a choisi de traiter. La largeur est le prix de la malédiction, et la meilleure largeur équilibre
deux effets (\cref{fig:m2:market-making-games:pnl}). Un marché plus large perd moins face aux informés,
qui n'interviennent que lorsque la somme se situe en dehors de celui-ci, mais il échange moins avec les
non-informés, qui paient à chaque fois la moitié de la largeur.

\begin{proposition}[Largeur face à la table]\label{prop:m2:market-making-games:width}
En l'absence d'intervenants informés, le profit espéré du teneur par tour vaut $\tfrac w2 e^{-w/\lambda}$
avec une patience $\lambda$, maximal pour $w = \lambda$. Les intervenants informés rendent chaque largeur plus coûteuse,
et les largeurs les plus étroites le sont davantage, puisqu'ils n'interviennent que contre des cotations placées du mauvais côté de
la valeur~; l'apprentissage à partir des transactions réduit la perte qu'ils infligent, et permet donc
un marché plus étroit.
\end{proposition}

\begin{proof}
Un intervenant non informé traite avec une probabilité $e^{-w/\lambda}$ et, le centre étant l'espérance de la valeur,
rapporte en moyenne $w/2$ au teneur~; le produit est maximal pour $w = \lambda$.
Le coût espéré pour le teneur d'un intervenant informé vaut $E[(V - a)^+ + (b - V)^+]$, qui diminue
à mesure que les cotations s'écartent~; un meilleur centre le réduit pour toute largeur.
\end{proof}

\begin{example}[La meilleure largeur]\label{ex:m2:market-making-games:best}
Sur 2\,000 parties simulées avec graine fixée, de 20 tours chacune, avec un cinquième des intervenants informés, le teneur qui met à jour
ses estimations gagne le plus, 16.4 par partie, à une largeur de 11.5 points, contre 14.1 à 8 points et
15.1 à 15 points. Un teneur qui ne met jamais à jour fait mieux avec une largeur plus grande, à 14 points, et ne gagne que 12.0.
Sans intervenants informés, la meilleure largeur serait la patience, soit 8 points, pour un gain d'environ
29 par partie (29.4 selon la formule, 28.6 en simulation).
\end{example}

\begin{omfigure}
\begin{tikzpicture}
\begin{axis}[width=11cm,height=5cm,
  xlabel={largeur du marché (points)},ylabel={P\&L du teneur par partie},
  tick label style={font=\small},label style={font=\small},
  xmin=0,xmax=30,ymin=-30,ymax=25,grid=major,grid style={gray!25},
  legend columns=4,legend style={font=\footnotesize,at={(0.5,-0.3)},anchor=north,draw=none,fill=none,/tikz/every even column/.append style={column sep=6pt}},
  table/col sep=comma]
\addplot[omDef,very thick] table[x=width,y=learn10]{figdata/markets-2/30-market-making-games/pnl.csv};
\addplot[omThm,very thick] table[x=width,y=learn20]{figdata/markets-2/30-market-making-games/pnl.csv};
\addplot[omProp,very thick] table[x=width,y=learn30]{figdata/markets-2/30-market-making-games/pnl.csv};
\addplot[omThm,thick,dashed] table[x=width,y=naive20]{figdata/markets-2/30-market-making-games/pnl.csv};
\addplot[gray,thin,forget plot] coordinates {(0,0) (30,0)};
\legend{10~\%\ d'informés,20~\%,30~\%,20~\%~: sans mise à jour}
\end{axis}
\end{tikzpicture}

\omcaption{Profit espéré du teneur par partie de 20 tours en fonction de la largeur de son marché,
pour des tables comptant 10, 20 et 30~\% d'intervenants informés, avec mise à jour selon la règle de Bayes, et pour 20~\%
sans mise à jour (tirets). Les marchés étroits se font cueillir~; les marchés larges ne sont pas traités~; la meilleure largeur
augmente, et le profit baisse, à mesure que la part des intervenants informés grandit. 600 parties simulées avec graine fixée par point.
Données~: le tutoriel du chapitre.}\label{fig:m2:market-making-games:pnl}
\end{omfigure}

\section{Les jeux d'entraînement que décrivent les entreprises}

Les entreprises publient rarement leurs jeux, mais certaines les décrivent. Le guide de Jane Street propose des marchés d'entraînement~:
le résultat d'un dé à vingt faces, le maximum de trois dés à six faces,
un dé avec la possibilité de le relancer une fois, la température dans une heure, et un million de fois le résultat d'un lancer de dé,
avec la question de savoir ce que l'on ressentirait si son cours acheteur était touché et que le dé tombait
sur un. Son stage de trading comprend des simulations de négociation, en équipes, sur des scénarios construits par ses traders,
ainsi que des cours qui vont des produits financiers au poker. Le fil conducteur est celui
du chapitre~: la valeur est incertaine, les autres joueurs ne savent pas les mêmes choses, et chaque
transaction est une information.

Les jeux apprennent aussi à gérer l'inventaire. Un teneur qui a vendu trois fois de suite est à découvert, et
si ses ventes ont été faites à des intervenants informés, la valeur se situe au-dessus du prix auquel il a vendu. Décaler ses deux cotations
vers le haut, ou les incliner, le protège, au prix de moins d'échanges d'un côté~; le moteur
de cartes peut être étendu pour le faire.

\section{Tutoriel~: le marché de la somme des cartes}\label{tut:m2:market-making-games:mm}

\begin{tutorial}
\textbf{Objectif.} Calculer la distribution exacte de la somme des cartes, faire jouer le marché contre une table
de robots informés et non informés, avec et sans \omterm{def:m2:market-making-games:bayes}{mise à jour bayésienne}, et trouver la
largeur qui maximise le profit espéré.
\textbf{État final~:} \cref{fig:m2:market-making-games:dist,fig:m2:market-making-games:game,fig:m2:market-making-games:pnl},
\cref{ex:m2:market-making-games:first,ex:m2:market-making-games:best} et les nombres du
problème du week-end.
\begin{tutsteps}
\item \textbf{Le jeu de cartes et la mise à jour}~: la distribution exacte de la somme, les paramètres de la table et la règle de Bayes
après une transaction ou une passe.
\omcode{code/firm/mmgame/firm_mmgame.py}{15}{49}{La distribution de la somme des cartes, la table et la mise à jour bayésienne.}\label{lst:m2:market-making-games:update}
\item \textbf{La partie}~: distribution des cartes et déroulement d'une partie.
\omcode{code/firm/mmgame/firm_mmgame.py}{52}{80}{Distribution et une partie.}\label{lst:m2:market-making-games:play}
\item \textbf{Les largeurs}~: le profit espéré de chaque largeur sur les mêmes donnes.
\omcode{code/firm/mmgame/firm_mmgame.py}{83}{95}{Profit espéré par largeur, sur des donnes communes.}\label{lst:m2:market-making-games:pnl}
\item \textbf{Exécuter} \texttt{mmgame\_demo.first\_trade()}, \texttt{mmgame\_demo.problem()}
et \texttt{fig\_mmgame.py}.
\end{tutsteps}
\textbf{Pour aller plus loin.} Laisser le teneur incliner ses cotations selon son inventaire, et vérifier
s'il gagne davantage~; puis donner à chaque robot l'une des cinq cartes au lieu de la somme, de sorte que
chaque négociateur soit en partie informé, comme autour d'une vraie table.
\end{tutorial}

\section{Construction~: le moteur du jeu de teneur de marché}\label{bld:m2:market-making-games:mmgame}

\begin{build}
\bfield{Rôle} La société miniature entraîne sa logique de cotation, et ses équipes, sur des jeux
dont elle peut calculer les réponses.
\bfield{Interface} \texttt{card\_sum\_distribution(k, copies, values)}~;
\texttt{Table(cards, rounds, informed, patience)}~; \texttt{mean(post)}~;
\texttt{update(post, action, bid, ask, table, width)}~; \texttt{deal(table, rng, prior)}~;
\texttt{play(width, table, value, draws, learn, prior)}~; \texttt{expected\_pnl(widths,
table, games, seed, learn)}. \bfield{Règles} Cinq cartes tirées d'un même jeu~; un trader par tour~; les traders informés n'échangent
que lorsque la somme est en dehors des cotations~; les traders non informés échangent avec la probabilité
$e^{-w/\lambda}$ d'un côté au hasard~; les cotations sont centrées sur la valeur espérée courante du teneur~;
chaque largeur fait face aux mêmes donnes tirées avec une même graine.
\bfield{Tests d'acceptation} \texttt{code/firm/mmgame/tests/}~: la distribution a une masse totale égale à un
d'espérance 35 et les bons extrêmes~; un achat relève et une vente abaisse l'estimation~; sans
trader informé, le profit suit la formule~; la \omterm{def:m2:market-making-games:bayes}{mise à jour bayésienne} fait mieux que l'absence de mise à jour~; les parties
sont reproductibles.
\bfield{Pour aller plus loin} Biais d'inventaire~; plusieurs teneurs de marché en concurrence sur la largeur~; traders disposant d'une information partielle (une carte chacun)~;
la fourchette à profit nul de Glosten--Milgrom~; une version en direct où des humains jouent
contre les robots.
\end{build}

\begin{omsources}
\item Jane Street, «~Probability and markets guide~»~; page du stage de trading.
\item L.~R.~Glosten and P.~R.~Milgrom, «~Bid, ask and transaction prices in a specialist
market with heterogeneously informed traders~», Journal of Financial Economics 14 (1985).
\end{omsources}

\section{Exercices}

\begin{exercise}[$\star$]\label{exo:m2:market-making-games:1}
Faire un marché sur le lancer d'un dé à six faces, et expliquer votre largeur.
\end{exercise}

\begin{exercise}[$\star$]\label{exo:m2:market-making-games:2}
On cote 200 à 400 sur le nombre de fenêtres et l'interviewer achète. Que fait-on
de son marché suivant, et pourquoi~?
\end{exercise}

\begin{exercise}[$\star$]\label{exo:m2:market-making-games:3}
Expliquer pourquoi une passe laisse l'estimation du teneur de marché à 35 au premier tour du jeu de cartes.
\end{exercise}

\begin{exercise}[$\star\star$]\label{exo:m2:market-making-games:4}
Sans trader informé et avec une patience de 8, déterminer la largeur qui maximise le
profit du teneur de marché par tour, et calculer ce qu'il gagne sur 20 tours.
\end{exercise}

\begin{exercise}[$\star\star$]\label{exo:m2:market-making-games:5}
Au premier tour, à 30 et 40, calculer la probabilité qu'un acheteur soit informé, puis
expliquer pourquoi elle est si nettement supérieure aux 20~\% de traders informés à la table.
\end{exercise}

\begin{exercise}[$\star\star$]\label{exo:m2:market-making-games:6}
Expliquer la \omterm{def:m2:market-making-games:curse}{malédiction du vainqueur} dans les termes d'un teneur de marché.
\end{exercise}

\begin{exercise}[$\star\star\star$]\label{exo:m2:market-making-games:7}
\emph{Programmation.} Avec \texttt{expected\_pnl}, trouver la meilleure largeur pour une table
avec 30~\% de traders informés sur la grille de 8 à 15 par demi-points, sur les mêmes 2\,000 donnes simulées.
\end{exercise}

\begin{exercise}[$\star\star\star$]\label{exo:m2:market-making-games:8}
\emph{Trouver l'erreur.} « Je cote serré pour trader davantage~: plus de transactions
rapportent plus de fourchette encaissée.~» Corriger ce raisonnement.
\end{exercise}

\section{Problème~: le marché de la somme des cartes}

\begin{problem}[{problème du week-end --- la largeur qui rapporte}]\label{pb:m2:market-making-games:1}
On fait un marché pendant 20 tours sur la somme de cinq cartes distribuées face cachée
dans un jeu de 52 cartes. À chaque tour, un trader se présente~: dans un cas sur cinq environ, il connaît la somme~; sinon,
un trader qui échange d'un côté au hasard, avec la probabilité $e^{-w/8}$ à la largeur $w$. On cote
autour de sa valeur espérée courante.

\medskip\noindent\textbf{Partie I --- Le contrat.}
\begin{enumerate}
\item Quelles sont l'espérance, l'écart type et l'étendue de la somme~?
\item Pourquoi les extrêmes sont-ils 6 et 64 plutôt que 5 et 65~?
\item Pour une largeur de 10, quelle est la probabilité qu'un trader non informé échange~?
\item Coter le premier marché à la largeur 10.
\item Combien gagne-t-on en moyenne par tour avec un trader non informé~?
\end{enumerate}

\medskip\noindent\textbf{Partie II --- Apprentissage.}
\begin{enumerate}[resume]
\item Après un achat à 40, quelle est la nouvelle valeur espérée~? Et après une vente à 30~?
\item Pourquoi une passe n'apprend-elle rien au premier tour, mais quelque chose plus tard~?
\item Dans le jeu d'exemple du chapitre, comment l'estimation évolue-t-elle, et pourquoi part-elle d'abord
dans le mauvais sens~?
\item Que rapporte la \omterm{def:m2:market-making-games:bayes}{mise à jour bayésienne}, aux meilleures largeurs, par rapport à l'absence de mise à jour~?
\item Pourquoi le teneur de marché qui ne met pas à jour préfère-t-il un marché plus large~?
\end{enumerate}

\medskip\noindent\textbf{Partie III --- La largeur.}
\begin{enumerate}[resume]
\item Trouver la meilleure largeur et le P\&L espéré par partie.
\item Que rapportent les largeurs 8 et 15~?
\item Sans trader informé, quelles seraient la meilleure largeur et le profit maximal~?
\item Comment la meilleure largeur varie-t-elle avec 10~\% et 30~\% de traders informés~?
\item Quelle est la probabilité qu'un trader non informé échange à la meilleure largeur~?
\end{enumerate}

\medskip\noindent\textbf{Partie IV --- Jugement.}
\begin{enumerate}[resume]
\item Comment sait-on, à une vraie table, quelle part des traders est informée~?
\item Comment faut-il ajuster ses cotations après trois ventes consécutives~?
\item Qu'apprend ce jeu sur la cotation de titres réels~?
\item Énoncer le \emph{résultat nommé}~: la largeur qui maximise le P\&L espéré face à cette
table, et ce P\&L.
\item En une phrase~: qu'est-ce que la largeur rapporte~?
\end{enumerate}
\end{problem}

\section{Questions d'entretien}

\begin{interviewq}[$\star$ \iqroles{trader}]\label{iq:m2:market-making-games:1}
Faites-moi un marché sur la somme de deux dés.
\end{interviewq}

\begin{interviewq}[$\star$ \iqroles{trader}]\label{iq:m2:market-making-games:2}
J'achète à votre cours vendeur. Que faites-vous maintenant~?
\end{interviewq}

\begin{interviewq}[$\star\star$ \iqroles{trader, chercheur}]\label{iq:m2:market-making-games:3}
Pourquoi une fourchette existe-t-elle même sans coûts~? Expliquez-le à l'aide d'un modèle.
\end{interviewq}

\begin{interviewq}[$\star\star$ \iqroles{trader}]\label{iq:m2:market-making-games:4}
Vous êtes vendeur net de trois contrats après trois ventes. Comment faites-vous évoluer votre marché~?
\end{interviewq}

\begin{interviewq}[$\star\star$ \iqroles{chercheur}]\label{iq:m2:market-making-games:5}
Comment estimeriez-vous la part de traders informés sur un marché à partir de ses transactions~?
\end{interviewq}

\begin{interviewq}[$\star\star\star$ \iqroles{développeur}]\label{iq:m2:market-making-games:6}
Concevriez-vous un serveur de jeu multijoueur de teneur de marché pour l'entraînement, avec des robots et des humains.
\end{interviewq}
